% Place abstract below.

\vspace{-0.3in}

Measuring the diffuse, highly-ionized baryonic content in galactic halos and the intergalactic medium through soft x-ray absorption spectroscopy of active galactic nuclei is a main scientific objective of the Lynx X-ray Observatory mission concept that can only be accomplished with a next-generation grating spectrometer. 
Realizing such an instrument using reflection grating technology requires thousands of custom blazed gratings that each perform with high diffraction efficiency to be manufactured and aligned to intercept radiation coming to a focus in a Wolter-I telescope. 
With Lynx performance requirements demanding significant improvements over Chandra and XMM-Newton in terms of both spectral sensitivity and spectral resolving power, previous approaches to blazed grating fabrication face limitations in their ability to meet these goals simultaneously. 

The aim of this thesis is to implement two recently-developed techniques in nanofabrication for this task, with an emphasis on beamline diffraction-efficiency testing for characterizing spectral sensitivity. 
In particular, thermally-activated selective topography equilibration (TASTE) is pursued as a means for fabricating a master grating with the key advantage that it enables blazed groove facets to be patterned in polymeric electron-beam resist over a non-parallel groove layout not limited by substrate crystal structure. 
%by electron-beam lithography sans dependence on substrate crystal structure that 
%leads to a degradation in resolving power. 
Additionally, substrate-conformal imprint lithography (SCIL) is studied as a method for mass manufacturing high-fidelity grating replicas in a silica sol-gel resist while avoiding many of the detriments associated with large-area patterning in other nanoimprint techniques. 

 % from a master template. %, a large-area nanoimprint technique,
Diffraction-efficiency testing of sub-micron grating prototypes coated with gold shows that TASTE is capable of meeting Lynx requirements for spectral sensitivity, with room for improvement at small groove periods, and that while SCIL offers a promising avenue for Lynx grating production, imprints suffer a small blaze-angle reduction due to resist shrinkage. 
With process development for TASTE established and equipment for SCIL recently installed at the Penn State Nanofabrication Laboratory, these findings will contribute to upcoming sounding-rocket experiments for x-ray spectroscopy. 
Accompanying this dissertation are appendices that outline physics fundamentals for x-ray spectral lines, x-ray optics, and diffraction gratings. 
% and more broadly, reflection grating solutions 
%grating imprints produced via SCIL are highly-effic

%SCIL yields imprints that, while highly-efficient, suffer a small blaze-angle reduction due to resist shrinkage. %is a promising avenue for Lynx grating production. 

%but a small level of resist shrinkage must be accounted for.  


%Additionally, substrate-conformal imprint lithography (SCIL)


%for the production of a master grating with a custom, blazed-groove layout and additionally, the mass manufacture of high-fidelity grating replicas.


%This motivates the pursuit of alternative fabrication methods for the production of a master grating with a custom, blazed-groove layout and additionally, the mass manufacture of high-fidelity grating replicas. 



%the recently-developed processes of thermally-activated selective topography equilibration (TASTE) and substrate-conformal imprint lithography (SCIL)



%\SI{40}{\percent} total diffraction efficiency and spectral resolving power greater than \num{5000} across the soft x-ray bandpass, previously-pursued nanofabrication techniques for blazed gratings face limitations in meeting these goals

